The roguelike game genre, known for its high replayability provided by
\gls{PCG}, often presents challenges regarding fine-grained control and the
creation of content with specific developer intentions. Manipulating complex and
rigid parameters in traditional PCG algorithms renders the design process
unintuitive. This work explores an approach to overcome this limitation by
proposing a system that uses \gls{LLM}s to generate maps from natural language
textual descriptions. Recognizing the inherent limitations of \gls{LLM}s in
spatial and geometric reasoning, the research developed a hybrid architecture
that decouples responsibilities: traditional deterministic algorithms manage the
structural topology and dungeon connectivity, while the \gls{LLM} acts as a
"creative agent"\ responsible for the semantic generation of narratives,
enemies, and items. To materialize these elements, a pipeline was implemented
integrating Structured Outputs and the \gls{RAG} technique over a curated vector
database of tiles, ensuring visual consistency and mitigating hallucinations. A
functional game prototype was developed to interpret the generated assets and
render playable levels. Validation, performed through Semantic Reconstruction
and Coherence metrics across three distinct models, demonstrated that smaller
models running locally offer performance comparable to large models. The results
confirm the technical feasibility of associating natural language with concrete
game assets, establishing a solid foundation for "text-to-map"\ creation systems
that preserve user intention.

\keywords{procedural content generation; large language models; roguelike games;
map generation.}